% Markets – Quantitative Analysis Summer Analyst — Citi
\documentclass[10pt,letterpaper]{article}

\usepackage[margin=0.75in]{geometry}
\usepackage[hidelinks]{hyperref}
\usepackage[skip=4pt]{parskip}

\pagestyle{empty}
\pagenumbering{gobble}
\setlength{\parindent}{0pt}

\begin{document}

\textbf{Rojae Mighty}\\
Stony Brook, NY \quad$\cdot$\quad
\href{mailto:rojae.mighty@stonybrook.edu}{rojae.mighty@stonybrook.edu}\\
\href{https://rojaemighty.com}{rojaemighty.com} \quad$\cdot$\quad
\href{https://www.linkedin.com/in/rojaemight/}{linkedin.com/in/rojaemight}

\vspace{8pt}
July 15, 2026

\vspace{8pt}
Hiring Manager\\
Citi\\
New York, NY

\vspace{8pt}
\textbf{Re: Markets -- Quantitative Analysis Summer Analyst}

\vspace{4pt}
Dear Hiring Manager,

I am excited to apply for the Markets -- Quantitative Analysis Summer Analyst position at Citi. I am an undergraduate physics student at Stony Brook University with research experience in computational modeling, statistical data analysis, and machine learning. Through my work at Brookhaven National Laboratory and the Center for Frontiers in Nuclear Science, I have developed the technical and analytical foundation to contribute to quantitative research in financial markets.

As a DOE Science Undergraduate Laboratory Intern at Brookhaven National Laboratory, I study electron--ion collisions using Monte Carlo event generation and computational analysis. I have implemented a spatially dependent parton distribution function and analyzed how nuclear environments affect particle behavior. This work requires me to translate physical questions into quantitative models, evaluate assumptions, analyze uncertainty, and interpret complex datasets using C++, Python, ROOT, and Unix/Linux tools.

At Stony Brook's Center for Frontiers in Nuclear Science, I investigate gluon saturation and geometric scaling at the Electron-Ion Collider. I also develop machine-learning methods for reconstructing deep-inelastic-scattering kinematics and build reproducible workflows for numerical simulations and Monte Carlo studies. These experiences have strengthened my ability to work systematically with data, validate models, and communicate technical results clearly.

In parallel, I am developing a market microstructure forecasting project using limit order book data. The project involves engineering features such as order-flow imbalance and spread dynamics, forecasting short-horizon mid-price returns and volatility, and applying walk-forward validation and transaction-cost-aware backtesting. This work has deepened my interest in quantitative finance and introduced me to the challenges of modeling noisy, time-dependent market data.

I am particularly drawn to Citi's Markets business because it combines rigorous quantitative analysis with fast-moving, real-world decision-making. I would be eager to bring my background in scientific modeling, computational research, and machine learning to Citi while continuing to develop my understanding of financial markets. I am confident that my curiosity, programming skills, and disciplined approach to analyzing complex systems would allow me to contribute effectively to the Quantitative Analysis team.

Thank you for considering my application. I would welcome the opportunity to discuss how my research experience, technical background, and interest in quantitative finance could contribute to Citi's Markets business.

\vspace{4pt}
Sincerely,\\[12pt]
Rojae Mighty

\end{document}
